% GENERATED FILE. Edit canonical lessons, then run npm run generate:books.
% canonical-source-hash: fnv1a64:3f0bbbb59036d580
% canonical-lessons: JA-C229-kaisha, JA-C229-kaigi, JA-C229-jimusho, JA-C229-tsutomeru, JA-C229-douryou, JA-R229-review-224-229

\chapter{At the Company}
\label{ch:ja-at-the-company}
\hlchaptermodality{229}

\begin{chapteropening}
\textbf{By the end of this chapter:} \emph{I can talk about companies, offices and colleagues.}

Complete the last lesson of chapter 229: I can talk about companies, offices and colleagues.
\end{chapteropening}

\section[\texorpdfstring{\ja{かいしゃ} (kaisha)}{かいしゃ (kaisha)}]{\ja{かいしゃ} (kaisha) — a company; a workplace}

\label{lesson:JA-C229-kaisha}



\begin{quote}
Before the new one: say the Japanese for to quit, then the Japanese for graduation.
\end{quote}

\subsection*{You'll want to know: \ja{かいしゃ}}
\textbf{\ja{かいしゃ}} — \emph{kaisha} — \textquotedblleft{}a company; a workplace\textquotedblright{}.

\begin{quote}
\textbf{\ja{ちちは} \ja{まいにち} \ja{かいしゃに} \ja{いきます。}} — \emph{chichi wa mainichi kaisha ni ikimasu} — \textquotedblleft{}My father goes to the company every day.\textquotedblright{}
\end{quote}

\begin{cousinweb}
\textbf{\ja{かいしゃ}} (\emph{kaisha}) is Sino-Japanese (kango): two Chinese characters read with their on'yomi, readings borrowed from Chinese. They are \emph{kai}, \textquotedblleft{}meet\textquotedblright{} (Middle Chinese \emph{hwajH}) and \emph{sha}, \textquotedblleft{}company, association\textquotedblright{} (\emph{dzyaeX}) (KANJIDIC2; Middle Chinese in Baxter's transcription).
\end{cousinweb}

\begin{grammarlens}[title={What you've built}]
One more word for everyday life at A2.
\end{grammarlens}

\subsection*{Your turn}
\begin{itemize}
\raggedright
  \item \emph{Say it:} \emph{kaisha}
  \item \emph{Say it:} \emph{kaisha}, once more
  \item \emph{Say it:} \emph{chichi wa mainichi kaisha ni ikimasu} — My father goes to the company every day.
  \item \emph{Recall:} say the Japanese for to quit, then the Japanese for graduation, then say \emph{kaisha} again
\end{itemize}

\begin{tcolorbox}[breakable,colback=teal!4,colframe=teal!35!black,title={Before you move on}]
What does \ja{かいしゃ} mean? (\textquotedblleft{}A company; a workplace\textquotedblright{}.) Say it once more.
\end{tcolorbox}
\section[\texorpdfstring{\ja{かいぎ} (kaigi)}{かいぎ (kaigi)}]{\ja{かいぎ} (kaigi) — a meeting; a conference}

\label{lesson:JA-C229-kaigi}



\begin{quote}
Before the new one: say the Japanese for graduation, then the Japanese for a company.
\end{quote}

\subsection*{You'll want to know: \ja{かいぎ}}
\textbf{\ja{かいぎ}} — \emph{kaigi} — \textquotedblleft{}a meeting; a conference\textquotedblright{}.

\begin{quote}
\textbf{\ja{かいしゃは} \ja{まいしゅう} \ja{かいぎが} \ja{あります。}} — \emph{kaisha wa maishū kaigi ga arimasu} — \textquotedblleft{}The company has a meeting every week.\textquotedblright{}
\end{quote}

\begin{cousinweb}
\textbf{\ja{かいぎ}} (\emph{kaigi}) is Sino-Japanese (kango): two Chinese characters read with their on'yomi, readings borrowed from Chinese. They are \emph{kai}, \textquotedblleft{}meet\textquotedblright{} (Middle Chinese \emph{hwajH}) and \emph{gi}, \textquotedblleft{}discuss\textquotedblright{} (\emph{ngjeH}) (KANJIDIC2; Middle Chinese in Baxter's transcription).
\end{cousinweb}

\begin{grammarlens}[title={What you've built}]
One more word for everyday life at A2.
\end{grammarlens}

\subsection*{Your turn}
\begin{itemize}
\raggedright
  \item \emph{Say it:} \emph{kaigi}
  \item \emph{Say it:} \emph{kaigi}, once more
  \item \emph{Say it:} \emph{kaisha wa maishū kaigi ga arimasu} — The company has a meeting every week.
  \item \emph{Recall:} say the Japanese for graduation, then the Japanese for a company, then say \emph{kaigi} again
\end{itemize}

\begin{tcolorbox}[breakable,colback=teal!4,colframe=teal!35!black,title={Before you move on}]
What does \ja{かいぎ} mean? (\textquotedblleft{}A meeting; a conference\textquotedblright{}.) Say it once more.
\end{tcolorbox}
\section[\texorpdfstring{\ja{じむしょ} (jimusho)}{じむしょ (jimusho)}]{\ja{じむしょ} (jimusho) — an office}

\label{lesson:JA-C229-jimusho}



\begin{quote}
Before the new one: say the Japanese for a company, then the Japanese for a meeting.
\end{quote}

\subsection*{You'll want to know: \ja{じむしょ}}
\textbf{\ja{じむしょ}} — \emph{jimusho} — \textquotedblleft{}an office\textquotedblright{}.

\begin{quote}
\textbf{\ja{ははは} \ja{えきの} \ja{ちかくの} \ja{じむしょで} \ja{はたらきます。}} — \emph{haha wa eki no chikaku no jimusho de hatarakimasu} — \textquotedblleft{}My mother works in an office near the station.\textquotedblright{}
\end{quote}

\begin{cousinweb}
\textbf{\ja{じむしょ}} (\emph{jimusho}) is Sino-Japanese (kango): three Chinese characters read with their on'yomi, readings borrowed from Chinese. They are \emph{ji}, \textquotedblleft{}affair\textquotedblright{} (Middle Chinese \emph{dzriH}), \emph{mu}, \textquotedblleft{}duty\textquotedblright{} (\emph{mjuH}) and \emph{sho}, \textquotedblleft{}place\textquotedblright{} (\emph{srjoX}) (KANJIDIC2; Middle Chinese in Baxter's transcription).
\end{cousinweb}

\begin{grammarlens}[title={What you've built}]
One more word for everyday life at A2.
\end{grammarlens}

\subsection*{Your turn}
\begin{itemize}
\raggedright
  \item \emph{Say it:} \emph{jimusho}
  \item \emph{Say it:} \emph{jimusho}, once more
  \item \emph{Say it:} \emph{haha wa eki no chikaku no jimusho de hatarakimasu} — My mother works in an office near the station.
  \item \emph{Recall:} say the Japanese for a company, then the Japanese for a meeting, then say \emph{jimusho} again
\end{itemize}

\begin{tcolorbox}[breakable,colback=teal!4,colframe=teal!35!black,title={Before you move on}]
What does \ja{じむしょ} mean? (\textquotedblleft{}An office\textquotedblright{}.) Say it once more.
\end{tcolorbox}
\section[\texorpdfstring{\ja{つとめる} (tsutomeru)}{つとめる (tsutomeru)}]{\ja{つとめる} (tsutomeru) — to be employed; to work for}

\label{lesson:JA-C229-tsutomeru}



\begin{quote}
Before the new one: say the Japanese for a meeting, then the Japanese for an office.
\end{quote}

\subsection*{You'll want to know: \ja{つとめる}}
\textbf{\ja{つとめる}} — \emph{tsutomeru} — \textquotedblleft{}to be employed; to work for\textquotedblright{}.

\begin{quote}
\textbf{\ja{あには} \ja{らいねん} \ja{ぎんこうに} \ja{つとめます。}} — \emph{ani wa rainen ginkō ni tsutomemasu} — \textquotedblleft{}Next year my elder brother will work for a bank.\textquotedblright{}
\end{quote}

Polite \textbf{\ja{つとめます}}, \emph{tsutomemasu}; an ichidan verb.

\begin{cousinweb}
\textbf{\ja{つとめる}} (\emph{tsutomeru}) is a native Japanese word (wago), inherited rather than borrowed from Chinese; its kanji was attached later as a kun'yomi (JMdict, KANJIDIC2).
\end{cousinweb}

\begin{grammarlens}[title={What you've built}]
One more word for everyday life at A2.
\end{grammarlens}

\subsection*{Your turn}
\begin{itemize}
\raggedright
  \item \emph{Say it:} \emph{tsutomeru}
  \item \emph{Say it:} \emph{tsutomeru}, once more
  \item \emph{Say it:} \emph{ani wa rainen ginkō ni tsutomemasu} — Next year my elder brother will work for a bank.
  \item \emph{Recall:} say the Japanese for a meeting, then the Japanese for an office, then say \emph{tsutomeru} again
\end{itemize}

\begin{tcolorbox}[breakable,colback=teal!4,colframe=teal!35!black,title={Before you move on}]
What does \ja{つとめる} mean? (\textquotedblleft{}To be employed; to work for\textquotedblright{}.) Say it once more.
\end{tcolorbox}
\section[\texorpdfstring{\ja{どうりょう} (dōryō)}{どうりょう (dōryō)}]{\ja{どうりょう} (dōryō) — a colleague; a co-worker}

\label{lesson:JA-C229-douryou}



\begin{quote}
Before the new one: say the Japanese for an office, then the Japanese for to be employed.
\end{quote}

\subsection*{You'll want to know: \ja{どうりょう}}
\textbf{\ja{どうりょう}} — \emph{dōryō} — \textquotedblleft{}a colleague; a co-worker\textquotedblright{}.

\begin{quote}
\textbf{\ja{ちちは} \ja{どうりょうと} \ja{よく} \ja{はなします。}} — \emph{chichi wa dōryō to yoku hanashimasu} — \textquotedblleft{}My father often talks with his colleagues.\textquotedblright{}
\end{quote}

\begin{cousinweb}
\textbf{\ja{どうりょう}} (\emph{dōryō}) is Sino-Japanese (kango): two Chinese characters read with their on'yomi, readings borrowed from Chinese. They are \emph{dō}, \textquotedblleft{}same\textquotedblright{} (Middle Chinese \emph{duwng}) and \emph{ryō}, \textquotedblleft{}companion, official\textquotedblright{} (\emph{lew}) (KANJIDIC2; Middle Chinese in Baxter's transcription).
\end{cousinweb}

\begin{grammarlens}[title={What you've built}]
One more word for everyday life at A2.
\end{grammarlens}

\subsection*{Your turn}
\begin{itemize}
\raggedright
  \item \emph{Say it:} \emph{dōryō}
  \item \emph{Say it:} \emph{dōryō}, once more
  \item \emph{Say it:} \emph{chichi wa dōryō to yoku hanashimasu} — My father often talks with his colleagues.
  \item \emph{Recall:} say the Japanese for an office, then the Japanese for to be employed, then say \emph{dōryō} again
\end{itemize}

\begin{tcolorbox}[breakable,colback=teal!4,colframe=teal!35!black,title={Before you move on}]
What does \ja{どうりょう} mean? (\textquotedblleft{}A colleague; a co-worker\textquotedblright{}.) Say it once more.
\end{tcolorbox}
\section[(dialogue)]{Work and school words, a review}

\label{lesson:JA-R229-review-224-229}



\begin{quote}
Say the Japanese for a student and a classroom.
\end{quote}

\subsection*{Your turn}
Say these aloud:

\begin{itemize}
\raggedright
  \item a student, a classroom, a class, to teach, to take lessons
  \item to memorize, practice, revision, to look up, a dictionary
  \item a question, to reply, an exam, to take an exam, a mark
  \item to make a mistake, to think over, to recall, to make sure, to compare
  \item to learn, to translate, to slack off, to quit, graduation
  \item a company, a meeting, an office, to be employed, a colleague
\end{itemize}

\begin{tcolorbox}[breakable,colback=teal!4,colframe=teal!35!black,title={Before you move on}]
A student? (\textbf{\ja{がくせい}}, \emph{gakusei}.) A classroom? (\textbf{\ja{きょうしつ}}, \emph{kyōshitsu}.) A class? (\textbf{\ja{じゅぎょう}}, \emph{jugyō}.) To teach? (\textbf{\ja{おしえる}}, \emph{oshieru}.) To take lessons? (\textbf{\ja{ならう}}, \emph{narau}.) To memorize? (\textbf{\ja{おぼえる}}, \emph{oboeru}.) Practice? (\textbf{\ja{れんしゅう}}, \emph{renshū}.) Revision? (\textbf{\ja{ふくしゅう}}, \emph{fukushū}.) To look up? (\textbf{\ja{しらべる}}, \emph{shiraberu}.) A dictionary? (\textbf{\ja{じしょ}}, \emph{jisho}.) A question? (\textbf{\ja{しつもん}}, \emph{shitsumon}.) To reply? (\textbf{\ja{こたえる}}, \emph{kotaeru}.) An exam? (\textbf{\ja{しけん}}, \emph{shiken}.) To take an exam? (\textbf{\ja{うける}}, \emph{ukeru}.) A mark? (\textbf{\ja{てん}}, \emph{ten}.) To make a mistake? (\textbf{\ja{まちがえる}}, \emph{machigaeru}.) To think over? (\textbf{\ja{かんがえる}}, \emph{kangaeru}.) To recall? (\textbf{\ja{おもいだす}}, \emph{omoidasu}.) To make sure? (\textbf{\ja{たしかめる}}, \emph{tashikameru}.) To compare? (\textbf{\ja{くらべる}}, \emph{kuraberu}.) To learn? (\textbf{\ja{まなぶ}}, \emph{manabu}.) To translate? (\textbf{\ja{やくす}}, \emph{yakusu}.) To slack off? (\textbf{\ja{なまける}}, \emph{namakeru}.) To quit? (\textbf{\ja{やめる}}, \emph{yameru}.) Graduation? (\textbf{\ja{そつぎょう}}, \emph{sotsugyō}.) A company? (\textbf{\ja{かいしゃ}}, \emph{kaisha}.) A meeting? (\textbf{\ja{かいぎ}}, \emph{kaigi}.) An office? (\textbf{\ja{じむしょ}}, \emph{jimusho}.) To be employed? (\textbf{\ja{つとめる}}, \emph{tsutomeru}.) A colleague? (\textbf{\ja{どうりょう}}, \emph{dōryō}.)
\end{tcolorbox}
